\documentclass[11pt]{ctexart}
\usepackage[a4paper,margin=1in]{geometry}
\usepackage{longtable,booktabs}
\title{Geant4-Agent 回归测试报告（中文）}
\author{自动化评测流程}
\date{2026-03-07\_123705}
\begin{document}
\maketitle
\section{统一 Pass/Fail 总表}
\begin{longtable}{p{0.06\linewidth}p{0.08\linewidth}p{0.07\linewidth}p{0.08\linewidth}p{0.07\linewidth}p{0.07\linewidth}p{0.07\linewidth}p{0.07\linewidth}p{0.09\linewidth}p{0.06\linewidth}p{0.21\linewidth}}
\toprule
ID & 领域 & 风格 & 参数完整性 & 期望初始 & 实际初始 & 期望最终 & 实际最终 & 缺参(初->终) & 结果 & 失败原因 \\
\midrule
MTM01 & medical & fuzzy & missing & 失败 & 失败 & 通过 & 通过 & 13->0 & 失败 & expected\_final\_mismatch:source.type \\
MTM02 & aerospace & precise & missing & 失败 & 失败 & 通过 & 通过 & 2->0 & 通过 & - \\
MTM03 & industrial\_ndt & precise & missing & 失败 & 失败 & 通过 & 通过 & 3->0 & 通过 & - \\
MTM04 & physics & precise & missing & 失败 & 失败 & 通过 & 通过 & 8->0 & 通过 & - \\
MTM05 & nuclear\_engineering & fuzzy & missing & 失败 & 失败 & 通过 & 通过 & 5->0 & 通过 & - \\
MTM06 & security\_screening & fuzzy & missing & 失败 & 失败 & 通过 & 失败 & 9->1 & 失败 & final\_complete\_mismatch \\
MTM07 & semiconductor & precise & missing & 失败 & 失败 & 通过 & 通过 & 13->0 & 失败 & expected\_final\_mismatch:source.type \\
MTM08 & education & fuzzy & missing & 失败 & 失败 & 通过 & 通过 & 11->0 & 通过 & - \\
MTM09 & environmental\_radiation & fuzzy & missing & 失败 & 失败 & 通过 & 通过 & 10->0 & 通过 & - \\
MTM10 & high\_energy\_physics & precise & missing & 失败 & 失败 & 通过 & 通过 & 10->0 & 通过 & - \\
MTM11 & materials\_update & precise & full & 通过 & 通过 & 通过 & 通过 & 0->0 & 通过 & - \\
MTM12 & overwrite\_reject & precise & full & 通过 & 通过 & 通过 & 通过 & 0->0 & 失败 & expected\_final\_mismatch:materials.selected\_materials.0, expected\_final\_mismatch:output.format \\
\bottomrule
\end{longtable}
\section{基准对比}
\begin{itemize}
\item pass\_rate=0.667
\item initial\_alignment\_rate=1.000
\item final\_alignment\_rate=0.917
\item avg\_turns\_per\_case=2.667
\item missing\_reduction\_rate=1.000
\item confirm\_apply\_success\_rate=1.000
\item rollback\_after\_confirm\_rate=0.000
\item internal\_leak\_turn\_rate=0.062
\item repeated\_assistant\_case\_rate=0.000
\item stale\_confirm\_case\_rate=0.000
\item delta\_pass\_rate=-0.283
\item regressions=无
\item improvements=无
\end{itemize}
\appendix
\section{附录：完整对话与内部处理链（自然语言）}
\subsection{案例 MTM01（medical）}
\textbf{对话：}
\begin{itemize}
\item 用户：???????????????????????????????????
\item 系统：你打算设计一个什么形状的几何体呢？如果是盒体的话，它的长度（x方向）大概是多少？
\item 用户：?? 300 mm x 300 mm x 300 mm ?????? G4\_WATER????? gamma??? 6 MeV??? (0,0,-400) mm??? +z????? QBBC??? root?
\item 系统：配置已完成。当前已确认的参数包括：物理列表的选择依据和来源，以及材料的选择依据和来源。
\end{itemize}
\textbf{内部链路（自然语言）：}
\begin{itemize}
\item 第 1 轮：动作=ask\_clarification，阶段=geometry，闭环状态=False。
\item 缺失字段：geometry.structure、geometry.params.module\_x、geometry.params.module\_y、geometry.params.module\_z、materials.volume\_material\_map、source.type、source.particle、source.energy、source.position、source.direction、physics.physics\_list、output.format、output.path。
\item LLM 使用=True；回退=无；后端=llm\_semantic\_frame+runtime\_semantic。
\item 第 2 轮：动作=finalize，阶段=materials，闭环状态=True。
\item 缺失字段：无。
\item LLM 使用=True；回退=无；后端=llm\_slot\_frame+runtime\_semantic。
\end{itemize}
\subsection{案例 MTM02（aerospace）}
\textbf{对话：}
\begin{itemize}
\item 用户：I need a first-pass proton irradiation setup for a silicon device; keep it simple for now.
\item 系统：Geometry: updated geometry structure, box size along x, box size along y. Materials: updated materials.selected\_materials, volume-to-material binding, material selection source. Source: updated source type, particle type, source energy. Physics: updated physics list, physics.backup\_physics\_list, physics selection reasons. Output: still needs output format, output path.
\item 用户：Use a silicon box 50 mm x 50 mm x 5 mm, beam source proton 100 MeV from (0,0,-100) mm along +z, physics FTFP\_BERT, output json.
\item 系统：You've already confirmed a geometry setup. This new request would change the following values:

*   Box size along x: from 1000.0 to 50.0
*   Box size along y: from 1000.0 to 50.0

Please reply **'confirm'** to apply these changes, or **'keep existing'** to keep the current values.
\item 用户：confirm
\item 系统：Your configuration is complete. The following items have been confirmed: geometry program, source type, and the box dimensions along the Z, Y, and X axes.
\end{itemize}
\textbf{内部链路（自然语言）：}
\begin{itemize}
\item 第 1 轮：动作=summarize\_progress，阶段=materials，闭环状态=False。
\item 缺失字段：output.format、output.path。
\item LLM 使用=True；回退=无；后端=llm\_slot\_frame+runtime\_semantic。
\item 第 2 轮：动作=confirm\_overwrite，阶段=source，闭环状态=False。
\item 缺失字段：无。
\item LLM 使用=True；回退=无；后端=llm\_slot\_frame+runtime\_semantic。
\item 第 3 轮：动作=finalize，阶段=physics，闭环状态=True。
\item 缺失字段：无。
\item LLM 使用=True；回退=无；后端=llm\_semantic\_frame+runtime\_semantic+confirmed\_pending\_overwrite。
\end{itemize}
\subsection{案例 MTM03（industrial\_ndt）}
\textbf{对话：}
\begin{itemize}
\item 用户：????? gamma ?????????????????????
\item 系统：目前我们已经确认了几何、材料和源的大部分设置。接下来还需要补充物理列表、输出格式和输出路径这三项信息。
\item 用户：????? 40 mm??? 80 mm ??????? G4\_Fe??? gamma??? 2 MeV??? (0,0,-200) mm??? +z????? Shielding??? csv?
\item 系统：我检测到这会覆盖已经确认的值。待修改内容：几何结构类型：single\_box -> single\_tubs；materials.selected\_materials：['G4\_Cu'] -> ['G4\_Fe']。如果要应用这些修改，请直接回复“确认”；如果要保留原值，请直接说“保持原值”。
\item 用户：确认
\item 系统：配置已完成。当前已确认的设置有：体积与材料绑定、半长、半径、源位置。
\end{itemize}
\textbf{内部链路（自然语言）：}
\begin{itemize}
\item 第 1 轮：动作=summarize\_progress，阶段=materials，闭环状态=False。
\item 缺失字段：physics.physics\_list、output.format、output.path。
\item LLM 使用=True；回退=无；后端=llm\_slot\_frame+runtime\_semantic。
\item 第 2 轮：动作=confirm\_overwrite，阶段=source，闭环状态=False。
\item 缺失字段：无。
\item LLM 使用=True；回退=无；后端=llm\_slot\_frame+runtime\_semantic。
\item 第 3 轮：动作=finalize，阶段=physics，闭环状态=True。
\item 缺失字段：无。
\item LLM 使用=True；回退=无；后端=llm\_semantic\_frame+runtime\_semantic+confirmed\_pending\_overwrite。
\end{itemize}
\subsection{案例 MTM04（physics）}
\textbf{对话：}
\begin{itemize}
\item 用户：Gamma point source 1 MeV at (0,0,-100) toward +z; please build the rest.
\item 系统：Geometry: still needs geometry structure, box size along x. Materials: still needs volume-to-material binding. Source: updated source type, particle type, source energy. Physics: still needs physics list. Output: still needs output format, output path.
\item 用户：Set copper box 1 m x 1 m x 1 m, physics FTFP\_BERT, output hdf5.
\item 系统：Your configuration is complete. The physics and material selections, along with their sources and reasons, have been confirmed.
\end{itemize}
\textbf{内部链路（自然语言）：}
\begin{itemize}
\item 第 1 轮：动作=summarize\_progress，阶段=geometry，闭环状态=False。
\item 缺失字段：geometry.structure、geometry.params.module\_x、geometry.params.module\_y、geometry.params.module\_z、materials.volume\_material\_map、physics.physics\_list、output.format、output.path。
\item LLM 使用=True；回退=无；后端=llm\_slot\_frame+runtime\_semantic。
\item 第 2 轮：动作=finalize，阶段=materials，闭环状态=True。
\item 缺失字段：无。
\item LLM 使用=True；回退=无；后端=llm\_slot\_frame+runtime\_semantic。
\end{itemize}
\subsection{案例 MTM05（nuclear\_engineering）}
\textbf{对话：}
\begin{itemize}
\item 用户：??????????????????????????
\item 系统：目前我们已经完成了几何结构类型和盒体尺寸的设置，以及源的能量、位置和方向。接下来还需要补充：体积与材料绑定、粒子类型、物理列表、输出格式和输出路径。
\item 用户：?? 200 mm x 200 mm x 200 mm ?????? G4\_Pb??? neutron??? 14 MeV??? (0,0,-50) mm??? +z????? QGSP\_BERT\_HP??? xml?
\item 系统：我注意到这将会覆盖已确认的参数。具体修改如下：盒体 x 方向尺寸从 1000.0 改为 200.0；盒体 y 方向尺寸从 1000.0 改为 200.0。如果您确认应用这些修改，请回复“确认”；如果您希望保持原值不变，请回复“保持原值”。
\item 用户：确认
\item 系统：配置已完成。当前已确认的参数包括：源位置、源能量、盒体 z 方向尺寸、盒体 y 方向尺寸。
\end{itemize}
\textbf{内部链路（自然语言）：}
\begin{itemize}
\item 第 1 轮：动作=summarize\_progress，阶段=materials，闭环状态=False。
\item 缺失字段：materials.volume\_material\_map、source.particle、physics.physics\_list、output.format、output.path。
\item LLM 使用=True；回退=无；后端=llm\_slot\_frame+runtime\_semantic。
\item 第 2 轮：动作=confirm\_overwrite，阶段=source，闭环状态=False。
\item 缺失字段：无。
\item LLM 使用=True；回退=无；后端=llm\_slot\_frame+runtime\_semantic。
\item 第 3 轮：动作=finalize，阶段=physics，闭环状态=True。
\item 缺失字段：无。
\item LLM 使用=True；回退=无；后端=llm\_semantic\_frame+runtime\_semantic+confirmed\_pending\_overwrite。
\end{itemize}
\subsection{案例 MTM06（security\_screening）}
\textbf{对话：}
\begin{itemize}
\item 用户：Need a quick cargo-screening toy setup with forward photons, but I have not fixed the target shape, material, source energy, or output format yet.
\item 系统：Here's a summary of your setup so far and what's still needed:

**Completed:**
*   **Source:** Particle type, position, and direction are set.

**Still Needed:**
*   **Geometry:** Structure and box dimensions (X, Y, Z).
*   **Materials:** Assignment of materials to your geometry volumes.
*   **Source:** Energy of the photons.
*   **Physics:** Selection of a physics list.
*   **Output:** Format and file path for your results.
\item 用户：Use steel cylinder target radius 60 mm half-length 120 mm, beam gamma 3 MeV from (0,0,-500) mm to +z, physics Shielding, output root.
\item 系统：You've requested a change that would overwrite two already confirmed settings:
*   Source type: from "point" to "beam".
*   Source position: from (0, 0, -100 mm) to (0, 0, -500 mm).

Please reply **'confirm'** to apply these changes, or **'keep existing'** to keep the current values.
\item 用户：confirm
\item 系统：Here's a summary of your setup progress:

**Updated this step:**
*   **Source:** Type and position have been updated as requested.

**Already confirmed:**
*   **Geometry:** Box dimensions (X, Y, Z).
*   **Physics:** The "Shielding" physics list.
*   **Output:** File path and ROOT format.

**Still needed:**
*   **Materials:** You must assign a material to your geometry volume.
\end{itemize}
\textbf{内部链路（自然语言）：}
\begin{itemize}
\item 第 1 轮：动作=summarize\_progress，阶段=geometry，闭环状态=False。
\item 缺失字段：geometry.structure、geometry.params.module\_x、geometry.params.module\_y、geometry.params.module\_z、materials.volume\_material\_map、source.energy、physics.physics\_list、output.format、output.path。
\item LLM 使用=True；回退=无；后端=llm\_slot\_frame+runtime\_semantic。
\item 第 2 轮：动作=confirm\_overwrite，阶段=materials，闭环状态=False。
\item 缺失字段：materials.volume\_material\_map。
\item LLM 使用=True；回退=无；后端=llm\_slot\_frame+runtime\_semantic。
\item 第 3 轮：动作=summarize\_progress，阶段=materials，闭环状态=False。
\item 缺失字段：materials.volume\_material\_map。
\item LLM 使用=True；回退=无；后端=llm\_semantic\_frame+runtime\_semantic+confirmed\_pending\_overwrite。
\end{itemize}
\subsection{案例 MTM07（semiconductor）}
\textbf{对话：}
\begin{itemize}
\item 用户：?????????????????????????
\item 系统：你打算模拟什么形状的几何体呢？比如一个简单的盒子？如果是盒子，它在 x 方向的尺寸大概是多少？
\item 用户：?? 20 mm x 20 mm x 1 mm ?????? G4\_Si????? neutron??? 2 MeV??? (0,0,0) mm????? QGSP\_BERT\_HP??? json?
\item 系统：配置已完成！当前已确认的设置有：物理列表选择依据、物理列表选择来源、材料选择依据和材料选择来源。
\end{itemize}
\textbf{内部链路（自然语言）：}
\begin{itemize}
\item 第 1 轮：动作=ask\_clarification，阶段=geometry，闭环状态=False。
\item 缺失字段：geometry.structure、geometry.params.module\_x、geometry.params.module\_y、geometry.params.module\_z、materials.volume\_material\_map、source.type、source.particle、source.energy、source.position、source.direction、physics.physics\_list、output.format、output.path。
\item LLM 使用=True；回退=无；后端=llm\_semantic\_frame+runtime\_semantic。
\item 第 2 轮：动作=finalize，阶段=materials，闭环状态=True。
\item 缺失字段：无。
\item LLM 使用=True；回退=无；后端=llm\_slot\_frame+runtime\_semantic。
\end{itemize}
\subsection{案例 MTM08（education）}
\textbf{对话：}
\begin{itemize}
\item 用户：For classroom demo, I want gamma attenuation in copper; you can ask only essential missing parts.
\item 系统：Here's your current setup summary:

**Completed:**
*   **Materials:** Copper is selected for the gamma attenuation demo.
*   **Source:** The particle type is set to gamma.
*   **Physics:** A suitable physics list has been selected for this simulation.

**Still needed to proceed:**
*   **Geometry:** The shape (structure) and its dimensions (box size along x, y, and z).
*   **Material Assignment:** You need to specify which geometry volume is made of copper.
*   **Source Details:** The source's type, energy, and position.
*   **Output:** The desired format and save location for your results.
\item 用户：Okay, copper box 100 mm x 100 mm x 100 mm, point gamma 1 MeV at (0,0,-200) mm, direction +z, physics FTFP\_BERT, output csv.
\item 系统：You're about to change a confirmed setting. The physics list will be updated from "Shielding" to "FTFP\_BERT". Please reply **'confirm'** to apply this change, or **'keep existing'** to keep the current value.
\item 用户：confirm
\item 系统：Your configuration is complete. The physics list, output path, output format, and source direction are all confirmed.
\end{itemize}
\textbf{内部链路（自然语言）：}
\begin{itemize}
\item 第 1 轮：动作=summarize\_progress，阶段=geometry，闭环状态=False。
\item 缺失字段：geometry.structure、geometry.params.module\_x、geometry.params.module\_y、geometry.params.module\_z、materials.volume\_material\_map、source.type、source.energy、source.position、source.direction、output.format、output.path。
\item LLM 使用=True；回退=无；后端=llm\_slot\_frame+runtime\_semantic。
\item 第 2 轮：动作=confirm\_overwrite，阶段=materials，闭环状态=False。
\item 缺失字段：无。
\item LLM 使用=True；回退=无；后端=llm\_slot\_frame+runtime\_semantic。
\item 第 3 轮：动作=finalize，阶段=source，闭环状态=True。
\item 缺失字段：无。
\item LLM 使用=True；回退=无；后端=llm\_semantic\_frame+runtime\_semantic+confirmed\_pending\_overwrite。
\end{itemize}
\subsection{案例 MTM09（environmental\_radiation）}
\textbf{对话：}
\begin{itemize}
\item 用户：?????????????? gamma ???????????
\item 系统：目前我们已经确定了源的设置：使用 gamma 粒子，并指定了来源和选择依据。

接下来还需要补充几何部分的参数，包括结构类型和尺寸，以及为体积指定材料。请告诉我这些信息，我们继续配置。
\item 用户：?? 500 mm x 500 mm x 500 mm ??????? G4\_AIR??? gamma??? 0.5 MeV??? (0,0,-300) mm??? +z????? QBBC??? root?
\item 系统：配置已完成。当前已确认的参数包括：物理列表的选择依据和来源，以及材料的选择依据和来源。
\end{itemize}
\textbf{内部链路（自然语言）：}
\begin{itemize}
\item 第 1 轮：动作=summarize\_progress，阶段=geometry，闭环状态=False。
\item 缺失字段：geometry.structure、geometry.params.module\_x、geometry.params.module\_y、geometry.params.module\_z、materials.volume\_material\_map、source.energy、source.position、physics.physics\_list、output.format、output.path。
\item LLM 使用=True；回退=无；后端=llm\_slot\_frame+runtime\_semantic。
\item 第 2 轮：动作=finalize，阶段=materials，闭环状态=True。
\item 缺失字段：无。
\item LLM 使用=True；回退=无；后端=llm\_slot\_frame+runtime\_semantic。
\end{itemize}
\subsection{案例 MTM10（high\_energy\_physics）}
\textbf{对话：}
\begin{itemize}
\item 用户：Set geometry as a sphere target; I have not decided radius or source details yet.
\item 系统：You've set up a sphere target geometry. Here's what's complete and what we still need:

**Completed:**
*   **Geometry:** Structure and root volume are set. The box dimensions around your sphere are confirmed.

**Still Needed:**
*   **Geometry:** The sphere's radius.
*   **Materials:** Which material to assign to your volume.
*   **Source:** The source type, particle type, and its energy.
*   **Physics:** The physics list to use.
*   **Output:** The output format and file path.
\item 用户：Sphere target radius 80 mm, lead material, electron point source 5 MeV at (0,0,-200) mm toward +z, physics FTFP\_BERT, output hdf5.
\item 系统：Your configuration is complete. The following items have been confirmed: the reasons for the physics list selection, the source of that selection, the reasons for the material choice, and the source of the material selection.
\end{itemize}
\textbf{内部链路（自然语言）：}
\begin{itemize}
\item 第 1 轮：动作=summarize\_progress，阶段=geometry，闭环状态=False。
\item 缺失字段：geometry.params.child\_rmax、materials.volume\_material\_map、source.type、source.particle、source.energy、source.position、source.direction、physics.physics\_list、output.format、output.path。
\item LLM 使用=True；回退=无；后端=llm\_slot\_frame+runtime\_semantic。
\item 第 2 轮：动作=finalize，阶段=materials，闭环状态=True。
\item 缺失字段：无。
\item LLM 使用=True；回退=无；后端=llm\_slot\_frame+runtime\_semantic。
\end{itemize}
\subsection{案例 MTM11（materials\_update）}
\textbf{对话：}
\begin{itemize}
\item 用户：Create a complete setup: 100 mm x 100 mm x 100 mm copper box, point gamma 1 MeV at (0,0,-200) mm toward +z, physics QBBC, output json.
\item 系统：Your configuration is complete. The following items have been confirmed: physics selection reasons, physics selection source, material selection reasons, and material selection source.
\item 用户：Change the target material to G4\_Al and change the output format to root.
\item 系统：You've requested changes that would overwrite already confirmed values.

Pending changes:
- Target material: from copper (G4\_Cu) to aluminum (G4\_Al)
- Output format: from JSON to ROOT

Please reply **"confirm"** to apply both changes, or **"keep existing"** to keep the current copper material and JSON format.
\item 用户：confirm
\item 系统：Configuration complete. The following items have been confirmed: volume-to-material binding, output path, output format, and selected materials.
\end{itemize}
\textbf{内部链路（自然语言）：}
\begin{itemize}
\item 第 1 轮：动作=finalize，阶段=materials，闭环状态=True。
\item 缺失字段：无。
\item LLM 使用=True；回退=无；后端=llm\_slot\_frame+runtime\_semantic。
\item 第 2 轮：动作=confirm\_overwrite，阶段=source，闭环状态=False。
\item 缺失字段：无。
\item LLM 使用=True；回退=无；后端=llm\_slot\_frame+runtime\_semantic。
\item 第 3 轮：动作=finalize，阶段=physics，闭环状态=True。
\item 缺失字段：无。
\item LLM 使用=True；回退=无；后端=llm\_semantic\_frame+runtime\_semantic+confirmed\_pending\_overwrite。
\end{itemize}
\subsection{案例 MTM12（overwrite\_reject）}
\textbf{对话：}
\begin{itemize}
\item 用户：??????????100 mm x 100 mm x 100 mm ?????? G4\_Cu??? gamma??? 1 MeV??? (0,0,-200) mm??? +z????? QBBC??? json?
\item 系统：配置已完成。当前已确认的参数包括：物理列表的选择依据和来源，以及材料的选择依据和来源。
\item 用户：????? G4\_Al????????? root?
\item 系统：我注意到这次修改会覆盖已经确认的参数。具体来说：

- 材料选择：从 `G4\_Cu` 改为 `G4\_Al`
- 输出格式：从 `json` 改为 `root`

如果您确认要应用这些修改，请直接回复“确认”；如果您希望保留原来的值，请直接说“保持原值”。
\item 用户：????
\item 系统：我注意到这次修改会覆盖已经确认的参数。具体来说：

- 材料选择：从 `G4\_Cu` 改为 `G4\_Al`
- 输出格式：从 `json` 改为 `root`

如果您确认要应用这些修改，请直接回复“确认”；如果您希望保留原来的值，请直接说“保持原值”。
\item 用户：确认
\item 系统：配置已完成。当前已确认的参数包括：体积与材料绑定、输出路径、输出格式以及材料选择（G4\_Al）。
\end{itemize}
\textbf{内部链路（自然语言）：}
\begin{itemize}
\item 第 1 轮：动作=finalize，阶段=materials，闭环状态=True。
\item 缺失字段：无。
\item LLM 使用=True；回退=无；后端=llm\_slot\_frame+runtime\_semantic。
\item 第 2 轮：动作=confirm\_overwrite，阶段=source，闭环状态=False。
\item 缺失字段：无。
\item LLM 使用=True；回退=无；后端=llm\_slot\_frame+runtime\_semantic。
\item 第 3 轮：动作=confirm\_overwrite，阶段=physics，闭环状态=False。
\item 缺失字段：无。
\item LLM 使用=True；回退=无；后端=llm\_slot\_frame+runtime\_semantic。
\item 第 4 轮：动作=finalize，阶段=output，闭环状态=True。
\item 缺失字段：无。
\item LLM 使用=True；回退=无；后端=llm\_semantic\_frame+runtime\_semantic+confirmed\_pending\_overwrite。
\end{itemize}
\end{document}